\documentclass[11pt]{article}

\usepackage[margin=0.85in]{geometry}
\usepackage{amsmath,amssymb}
\usepackage{booktabs}
\usepackage{array}
\usepackage{enumitem}
\usepackage{fancyhdr}
\usepackage{listings}
\usepackage{xcolor}
\usepackage{microtype}
\usepackage[hidelinks]{hyperref}

\newcommand{\assignmentnumber}{8}
\newcommand{\stagedue}{Monday, 14 September 2026, 5:00 PM IST}
\newcommand{\finaldue}{Monday, 21 September 2026, 11:59 PM IST}

\pagestyle{fancy}
\fancyhf{}
\lhead{DA3410}
\rhead{Assignment \assignmentnumber}
\cfoot{\thepage}
\renewcommand{\headrulewidth}{0.4pt}

\setlength{\parindent}{0pt}
\setlength{\parskip}{0.45em}
\setlist[itemize]{leftmargin=1.5em,itemsep=0.2em,topsep=0.25em}
\setlist[enumerate]{leftmargin=1.8em,itemsep=0.25em,topsep=0.25em}

\lstset{
    basicstyle=\ttfamily\small,
    frame=single,
    columns=fullflexible,
    breaklines=true,
    showstringspaces=false,
    keywordstyle=\bfseries,
    commentstyle=\itshape
}

\begin{document}

\begin{center}
    {\Large \textbf{DA3410: Assignment \assignmentnumber}}
\end{center}


\section{Objectives}
\label{sec:objectives}
 
\begin{itemize}
  \item Use the PatchCamelyon dataset in this lab. The dataset contains about
        327680 color patches. Each patch is 96 by 96 pixels. Each patch has a
        binary label. The label shows if the center region contains tumor tissue.
        The two classes are balanced.
  \item Reproduce a set of published results from the assigned paper on this data.
  \item Design controlled experiments and change one factor at a time.
  \item Measure the results, plot them, and compare them with the results in the
        papers.
  \item Explain the results in clear and simple terms. State where your results
        agree with the papers and where they differ.
\end{itemize}



\section{Deep Double Descent: Where Bigger Models and More Data Hurt}
\label{sec:double-descent}

\subsection*{Summary of Paper}
The classical bias-variance trade-off states that test error first goes down and
then goes up as a model becomes larger. This paper shows that modern networks do
not follow this rule. As the model becomes larger, the test error first goes down,
then goes up, and then goes down again. The peak in test error occurs at the
interpolation threshold. This is the point where the model is just large enough to
fit the training data with almost zero training error. The authors call this
behavior double descent.

The CNN experiments use a family of 5-layer convolutional networks. Each
network has four convolutional layers with widths $[k, 2k, 4k, 8k]$ and one fully
connected layer.

\subsection*{Reproducible Results on PatchCamelyon}
\begin{itemize}
  \item \textbf{Model-wise double descent:} Sweep the width $k$. Plot the test
        error against the model size. The plot shows the U-shape, the peak at the
        interpolation threshold, and the second descent.
  \item \textbf{Effect of label noise:} Repeat the sweep at $0\%$, $10\%$, and
        $20\%$ label noise. The peak grows and moves to larger models as the
        noise increases.
  \item \textbf{Epoch-wise double descent:} Fix one large CNN. Train it for many
        epochs. Plot the test error against the epoch number.
  \item \textbf{Interpolation and the peak:} Plot the training error on the same
        axis as the model-size sweep. The test peak appears where the training
        error first reaches near zero.
  \item \textbf{Effect of data augmentation:} Run the sweep with and without
        augmentation. Augmentation moves the interpolation threshold and the peak
        to larger models.
  \item \textbf{Effect of the optimizer:} Run the sweep with SGD and with Adam.
        Compare the position and the height of the peak.
  \item \textbf{Sample-wise behavior:} Subsample the PatchCamelyon training set.
        Find a range of model sizes where more data does not lower the test error
        at convergence.
  \item \textbf{Ensembling:} Train several models and combine their
        predictions.
\end{itemize}

\section{ConvNeXt V2: Co-designing and Scaling ConvNets with Masked Autoencoders}
\label{sec:convnextv2}

\subsection*{Summary of Paper}
This paper makes pure convolutional networks work well with masked self-supervised
learning. It has two main parts. The first part is the Fully Convolutional Masked
Autoencoder (FCMAE). The method hides about $60\%$ of the input patches. It
processes only the visible pixels with sparse convolutions. This step stops
information from the hidden regions from reaching the model. A small ConvNeXt
decoder then rebuilds the hidden patches. The loss is the mean squared error on
the hidden patches only.

\subsection*{Reproducible Results on PatchCamelyon}
\begin{itemize}
  \item \textbf{FCMAE pretraining:} Pretrain the model on unlabeled
        PatchCamelyon with masked reconstruction. Then fine-tune it on labeled
        PatchCamelyon. Show that the pretrained model reaches higher accuracy
        than a model that is trained from scratch. Change the labeled fraction to
        $1\%$, $10\%$, and $100\%$. Show that the gain from pretraining grows as
        the number of labels drops.
  \item \textbf{Effect of Global Response Normalization (GRN):} Compare the V1 block with the V2 block that has
        GRN. Also compare GRN with Batch, Layer, and Local Response Normalization
        in the same place.
  \item \textbf{Feature collapse analysis:} Show the activation map of each
        channel. Compute the mean pairwise cosine distance between channels for
        each layer. Show the collapse for FCMAE without GRN.
  \item \textbf{Masking ratio:} Change the masking ratio from $0.1$ to $0.9$.
  \item \textbf{Decoder design:} Change the depth and the width of the single
        block decoder.
  \item \textbf{Information leakage:} Compare correct masked encoding with plain
        dense convolution over the masked input.
\end{itemize}


\section{Tent: Fully Test-Time Adaptation by Entropy Minimization}
\label{sec:tent}

\subsection*{Summary of Paper}
This paper adapts a trained model to a new test distribution at test time. The
method uses only the unlabeled target data and the model itself. It does not use
the source data. It does not use labels. It does not change the training process.
The method is called Tent. Tent lowers the Shannon entropy of the model
predictions on each test batch.

\subsection*{Reproducible results on PatchCamelyon}
\begin{itemize}
  \item \textbf{Test-time BN update:} Compute the Batch Normalization statistics
        again on the target data. Use these statistics in place of the source
        statistics.
  \item \textbf{Tent result:} Freeze all weights except the BN affine
        parameters. Compute the statistics again for each batch. Minimize the
        prediction entropy.
  \item \textbf{Entropy and error:} Plot the entropy of each prediction against
        the error rate on corrupted PatchCamelyon.
  \item \textbf{Severity, entropy, and loss:} Show that the entropy and the loss
        both rise as the corruption severity rises.
  \item \textbf{Effect across corruption types:} Apply several corruption types.
        Show that Tent improves most of them.
  \item \textbf{Choice of parameters to update:} Compare three options. Update
        the affine parameters only. Update the last layer. Update the full model.
        Show that the full update makes the result worse.
  \item \textbf{Online and offline adaptation:} Compare one gradient step for
        each batch with several epochs of adaptation.
\end{itemize}

\end{document}